\chapter{Verificación posterior a la instalación}
\label{cap:verificacion}

Recorra esta lista al terminar cualquier instalación y anote el resultado en el registro.

\begin{xltabular}{\linewidth}{@{}c L L@{}}
	\toprule
	\textbf{\#} & \textbf{Comprobación} & \textbf{Resultado esperado} \\
	\midrule
	\endhead
	1 & Todos los servicios en marcha (Docker: \texttt{docker compose ps -a}; Linux:
	\texttt{systemctl status}; Windows: \texttt{Get-Service}) & En marcha; \texttt{setup}
	terminado con código 0 \\
	2 & Desde una PC cliente, abrir \texttt{https://<DIGAIPE\_HOST>} & Página de inicio sin aviso
	de certificado \\
	3 & Desde una PC cliente, abrir \texttt{http://<DIGAIPE\_HOST>} & Redirige a
	\texttt{https://} \\
	4 & Desde una PC cliente, probar los puertos internos: \texttt{Test-NetConnection
	<DIGAIPE\_HOST> -Port 3306} y \texttt{-Port 3000} & \texttt{TcpTestSucceeded : False} en
	ambos \\
	5 & Iniciar sesión con \texttt{BOOTSTRAP\_ADMIN\_EMAIL} & Entra a DIGAIPE \\
	6 & Administración → Auditoría & Aparece el inicio de sesión del punto 5, con la IP de la PC
	cliente (no \texttt{127.0.0.1} ni la IP del servidor) \\
	7 & Crear un usuario de prueba con un rol, iniciar sesión con él y comprobar que solo ve lo
	que permite su rol & Las páginas sin permiso responden \emph{acceso denegado} y quedan en la
	auditoría como \texttt{route.access\_denied} \\
	8 & Cargar un lote de admisiones de prueba y generar una planilla PDF & La carga y el PDF
	funcionan \\
	9 & Primer respaldo (registros de \texttt{backup-worker}, \texttt{digaipe-respaldo} o
	\texttt{respaldo.log}) & \texttt{Backup finished.}; en Auditoría, un evento
	\texttt{backup.completed} \\
	10 & Verificar la cadena de auditoría (capítulo~\ref{cap:auditoria}) & Verificación correcta,
	sin triggers faltantes \\
	11 & Reiniciar el servidor y repetir 1, 2 y 5 & Todo arranca solo \\
	12 & Borrar el usuario y los datos de prueba de los puntos 7 y 8, si no deben quedar & — \\
	\bottomrule
\end{xltabular}

\begin{nota}[Punto 6]
	Si la auditoría muestra siempre la misma IP (la del servidor o \texttt{127.0.0.1}), la
	aplicación no está leyendo la cabecera \texttt{X-Forwarded-For}: revise
	\texttt{ADDRESS\_HEADER} y \texttt{XFF\_DEPTH} (capítulo~\ref{cap:problemas}). Sin la IP real,
	el límite de intentos de inicio de sesión bloquearía a todos los usuarios a la vez.
\end{nota}

\begin{aviso}[Después de verificar]
	\begin{itemize}
		\item Cambie o retire \texttt{BOOTSTRAP\_ADMIN\_PASSWORD} del archivo de configuración si su
		      política lo exige: el administrador ya existe y no se vuelve a crear.
		\item Haga la primera copia del respaldo a un medio externo y el primer anclaje de
		      auditoría.
	\end{itemize}
\end{aviso}
